\begin{proposition}[The motivic groups of the diagonal data]\hypertarget{tag-001C}{}\label{ari:C.7.3}
\StatementRecord{ari:C.7.3}{proved}{001C}{\SeriesRef{ari:eq:cp-MTdim}}

Let $P$ be nontorsion and let $R$ be any algebraic lift above
$P$ in $G_P$. Put $M_R=[\Z\to G_P]$, $1\mapsto R$.
Its connected motivic group is, after a choice of Levi subgroup,
the semidirect product of $G_{\mathrm{mot}}(E)$ and the
three-dimensional Heisenberg group attached to $H_1(E)$ and
its principal polarization. Its dimension is $5$ for CM and
$7$ without CM.

For a non-real embedding put
$\mathcal M=M_R\oplus\overline{M_R}$.
If $E$ and $\bar E$ are not isogenous, its dimension is $13$.
If they are isogenous, choose an isogeny
$\phi:\bar E\to E$ of degree $d$ and write
$Q=\phi(\bar P)$. Let $r$ be the dimension over
$F=\End(E)\otimes\Q$ of the span of $P,Q$, modulo torsion.
Then
\[
\begin{array}{c|c|c|c}
\text{case}&\dim B&\dim Z&\dim G_{\mathrm{mot}}(\mathcal M)\\ \hline
r=2,\ E\text{ CM}&2&2&8\\
r=2,\ E\text{ not CM}&2&2&10\\
r=1,\ E\text{ CM}&1&1+e&5+e\\
r=1,\ E\text{ not CM}&1&1+e&7+e.
\end{array}
\]
Here $e$ is either zero or one: it is zero precisely when the
algebraic lift in the one-dimensional quotient of the two
central tori by the bracket torus is torsion. This criterion is
independent of the choices made modulo isogeny. In the real
case we use the single motive $M_R$; no extra quotient torus
is necessary.
\end{proposition}


\begin{proof}
\begin{enumerate}[label=\arabic*.,leftmargin=*]
\item All curves, polarizations, points, and lifts are algebraic over
$\overline\Q$, as required by Construction 3.4 and Lemma 3.5.
For $M_R$ the smallest abelian subvariety containing $(P,P)$ is
the diagonal $B=E$, since $P$ is nontorsion. If $\psi$ is the
polarization pairing, the bracket on this diagonal is
\[
 [(x,x),(y,y)]=\psi(x,y)-\psi(y,x)=2\psi(x,y).
\]
It is nonzero and spans $\Q(1)$. Hence $Z'=Z=\mathbb G_m$;
there is no deficient or Ribet exception, including for CM.
The pure group is $\mathrm{GL}_2$ without CM and
$\operatorname{Res}_{F/\Q}\mathbb G_m$ for the imaginary
quadratic CM field. Their dimensions are four and two.
This proves the single-motive assertion, including its Lie
bracket and Levi action: the action on the centre is the
polarization multiplier.

\item For the direct sum only the two diagonal central coordinates
occur. If the two elliptic curves are not isogenous, they
cannot be CM conjugates. The pure group is
$\{(g_1,g_2)\in\mathrm{GL}_2^2:\det g_1=\det g_2\}$,
of dimension seven. The point data generate $E\times\bar E$,
and the two brackets are independent, giving $7+4+2=13$.

\item If the curves are isogenous, transport the second polarized
factor through $\phi$. Its polarization pairing becomes
$\psi/d$. The smallest abelian subvariety has dimension $r$.
For $r=2$, independent variation in its two factors makes the
two diagonal brackets span both central coordinates; thus
$Z'=Z=\mathbb G_m^2$. The unipotent group is the product
of two Heisenberg groups, with the same pure group acting
on both standard factors.
For $r=1$, after multiplying the points by a common integer
we may write $Q=aP$, $a\in F^\times$. The bracket image is
the one-dimensional line with proportional central coordinates
\begin{equation}\hypertarget{tag-003H}{}\label{ari:eq:cp-positive-centre}
 (1,c),\qquad c=\frac{a a^\dagger}{d}>0,
\end{equation}
where $\dagger$ is the Rosati involution. On an elliptic curve
$aa^\dagger$ is a positive rational number: a square in the
non-CM case and the positive field norm in the CM case.
Construction 3.4(3) adds a further central direction exactly
when the image of the algebraic lift in the quotient torus
is nontorsion. This is the integer $e$ above.
The unipotent Lie algebra is the Heisenberg algebra, with one
additional central Tate line when $e=1$. Formula
\eqref{ari:eq:cp-MTdim} proves every entry of the table. This
description computes the group even when the lift contributes
a Kummer factor; specifying only the span of $P,Q$ would
not determine that last factor.
\end{enumerate}
\end{proof}